\documentclass[10pt]{article}
\usepackage[T1]{fontenc}
\usepackage[utf8]{inputenc}
\usepackage[french]{babel}
\usepackage{amsmath,amssymb,mathrsfs,mathpazo}
\usepackage{geometry}
\geometry{a4paper,margin=2cm}
\usepackage{graphicx,array,tabularx,tabularray,colortbl,multirow,multicol,booktabs,enumitem,pifont,ulem,fancybox,diagbox,cancel}
\usepackage{pstricks,pstricks-add,pst-plot,pst-tree,pst-node,pst-text,pst-func,pst-math,pst-bspline,pst-3dplot,pst-eucl,pst-solides3d}
\let\pscalculate\undefined
\usepackage{pst-bezier}
\usepackage{tikz,pgfplots,tkz-tab}
\usetikzlibrary{arrows,arrows.meta,patterns,calc,positioning,trees,decorations.markings,shapes.geometric,shapes.misc,shapes.arrows,intersections,decorations.pathmorphing,decorations.pathreplacing,backgrounds,shadows,babel}
\pgfplotsset{compat=1.18}
\usepgfplotslibrary{fillbetween}
\usepackage[np]{numprint}
\usepackage{listings,makecell,pas-tableur,twoopt}
\let\barre\undefined
\usepackage[pstricks]{tablvar}
\definecolor{darkgreen}{RGB}{0,100,0}
\providecommand{\red}{\color{red}}
\providecommand{\blue}{\color{blue}}
\providecommand{\green}{\color{green}}
\providecommand{\black}{\color{black}}
\providecommand{\white}{\color{white}}
\providecommand{\gray}{\color{gray}}
\providecommand{\cyan}{\color{cyan}}
\providecommand{\magenta}{\color{magenta}}
\providecommand{\yellow}{\color{yellow}}
\providecommand{\violet}{\color{violet}}
\providecommand{\orange}{\color{orange}}
\providecommand{\euro}{\texteuro}
\providecommand{\EUR}{\texteuro}
\providecommand{\ieme}{\textsuperscript{e}}
\providecommand{\degres}{\ensuremath{^\circ}}
\providecommand{\degree}{\ensuremath{^\circ}}
\providecommand{\vv}[1]{\overrightarrow{#1}}
\providecommand{\Coord}[1]{#1}
\providecommand{\decosix}{}
\providecommand{\starredbullet}{\ensuremath{\bullet}}
\pagestyle{empty}
\setlength{\parindent}{0pt}
\usepackage{adjustbox}\DeclareUnicodeCharacter{2236}{:}\DeclareUnicodeCharacter{2190}{\ensuremath{\leftarrow}}
\providecommand{\R}{\mathbb{R}}
\providecommand{\N}{\mathbb{N}}
\providecommand{\Z}{\mathbb{Z}}
\providecommand{\Q}{\mathbb{Q}}
\providecommand{\D}{\mathbb{D}}
\providecommand{\C}{\mathbb{C}}
\providecommand{\e}{{\mathrm{e}}}
\providecommand{\vect}[1]{\overrightarrow{#1}}
\providecommand{\vectt}[1]{\overrightarrow{\mathrm{#1}}}
\providecommand{\vvt}[1]{\overrightarrow{\mathrm{#1}}}
\providecommand{\np}[2][]{{#2~#1}}
\providecommand{\vv}[1]{{\overrightarrow{#1}}}
\renewcommand{\d}{{\mathrm{d}}}
\providecommand{\degres}{{^\circ}}
\providecommand{\degre}{{^\circ}}
\providecommand{\blue}{\color{blue}}
\providecommand{\red}{\color{red}}
\providecommand{\enskip}{\quad}
\providecommand{\llbracket}{\lbrack\!\lbrack}
\providecommand{\rrbracket}{\rbrack\!\rbrack}
\providecommand{\euro}{{\text{€}}}
\providecommand{\no}{{\text{n°}}}
\providecommand{\interval}[2]{{[#1;#2]}}
\providecommand{\intervFO}[2]{{[#1;#2[}}
\providecommand{\dotfill}{\ldots}

\providecommand{\ds}{\displaystyle}
\providecommand{\textding}[1]{#1}
\providecommand{\Oij}{\(\mathrm{O};\vec{\imath},\vec{\jmath}\)}
\providecommand{\Oijk}{\(\mathrm{O};\vec{\imath},\vec{\jmath},\vec{k}\)}
\providecommand{\Ouv}{\(\mathrm{O};\vec{u},\vec{v}\)}
\providecommand{\intervFF}[2]{\ensuremath{[#1;#2]}}
\renewcommand{\euro}{\ifmmode\text{\texteuro}\else\texteuro\fi}
\DeclareUnicodeCharacter{00B0}{\ifmmode^{\circ}\else\textdegree\fi}
\DeclareUnicodeCharacter{20AC}{\ifmmode\text{\texteuro}\else\texteuro\fi}
% Énoncés : Éducation nationale. Transcription et corrigés : APMEP, Denis Vergès et contributeurs. Découpage et classement : CoopMaths, Sébastien Lozano, Éric Elter et contributeurs. Adaptation de la présentation : Maths Explorer.
% CC BY-SA : https://coopmaths.fr/www/about/
% https://coopmaths.fr/alea/static/bac/2023/tex/bac_2023_09_sujet2_metropole_2.tex
\begin{document}
\noindent\textit{APMEP / CoopMaths. CC BY-SA. Présentation adaptée par Maths Explorer.}\par\medskip
\medskip

On considère la fonction $f$ définie sur l'intervalle $] 0~;~+\infty[$ par 
\[f(x) = (2-\ln x) \times \ln x,\]

où $\ln$ désigne la fonction logarithme népérien.

On admet que la fonction $f$ est deux fois dérivable sur $]0~;~+\infty[$.

On note $C$ la courbe représentative de la fonction $f$ dans un repère orthogonal et $C'$ la courbe représentative de la fonction $f'$, fonction dérivée de la fonction $f$.

La \textbf{courbe} $\boldsymbol{C}'$ est donnée ci-dessous ainsi que son unique tangente horizontale (T).

\begin{center}
\begin{tikzpicture}[x=6.35mm,y=35mm]






	\draw [color = gray!75, line width = 0.5pt, xstep=1, ystep = 0.2,dashed] (0,-0.39) grid (17.5,2.19);
\draw[->] (-0.5,0) -- (17.5,0);
\foreach \x in {1,...,17}
\draw[shift={(\x,0)},color=black] (0pt,2pt) -- (0pt,-2pt) node[below, fill = white] {\footnotesize $\np{\x}$};
\draw[->] (0,-0.39) -- (0,2.19);
\foreach \y in {-0.2,0.2,0.4,0.6,0.8,1,1.2,1.4,1.6,1.8,2}
\draw[shift={(0,\y)},color=black] (2pt,0pt) -- (-2pt,0pt) node[left, fill = white] {\footnotesize $\np{\y}$};
\draw[color=black] (-2pt,-2pt) node[below left] {O};
\clip (0,2.2) rectangle (17.5,-0.4);
\draw[line width=1pt,color=red,smooth,samples=100,domain= 0.8:17.5 ] plot(\x,{2*(1 - ln(\x))/\x} );
\node[red] at(1.65,1.3){$C'$};
\draw[teal,line width=1pt] (-0.5,-0.2707)--(17.5,-0.2707) ;
\draw[teal,line width=1pt,dashed] (7.3891,-0.2707)--(7.3891,0) ;
\node[teal] at(0.8,-0.35)  {(T)};
\draw[fill=white,draw=gray,line width=1pt] (3,1.4)rectangle(16,2); \node[text width=8.2cm,align=center] at (9.75,1.7){On rappelle que \textbf{cette courbe $C'$ est la courbe représentative de la fonction dérivée $f'$}} ;
\end{tikzpicture}
\end{center}


\begin{enumerate}
	\item Par lecture graphique, avec la précision que permet le tracé ci-dessus, donner :

	\begin{enumerate}
		\item le coefficient directeur de la tangente à $C$ au point d'abscisse 1.

		\item 	le plus grand intervalle sur lequel la fonction $f$ est convexe.
	\end{enumerate}

	\item \begin{enumerate}
		\item Calculer la limite de la fonction $f$ en $+\infty$.

		\item Calculer $\lim\limits_{x \to 0} f(x)$. Interpréter graphiquement ce résultat.
	\end{enumerate}

	\item Montrer que la courbe $C$ coupe l'axe des abscisses en deux points exactement dont on précisera les coordonnées.

	\item \begin{enumerate}
		\item Montrer que pour tout réel $x$ appartenant à $]0~;~+\infty[, \quad f'(x) = \dfrac{2(1-\ln x)}{x}$.

		\item En déduire, en justifiant, le tableau de variations de la fonction $f$ sur $]0~;~+\infty[$.
	\end{enumerate}

	\item On note $f''$ la dérivée seconde de $f$ et on admet que pour tout réel $x$ appartenant à $]0~;~+\infty[$, $f''(x) = \dfrac{2(\ln x-2)}{x^{2}}$.

Déterminer par le calcul le plus grand intervalle sur lequel la fonction $f$ est convexe et préciser les coordonnées du point d'inflexion de la courbe $C$.
\end{enumerate}

\vspace{1cm}

\end{document}
